% Section 11.7, from lectures/L11/README.md.

\section{Summary}
\label{c11:sec:review}

\needspace{6\baselineskip}
\subsection*{What you should be able to do now}
\begin{itemize}
  \item Draw the controller's block diagram and say what every block is responsible for, and just as
    importantly, what it is not responsible for.
  \item Write the top level's entity with the right port order and the right types.
  \item Explain why the architecture is empty now and which chapter fills which part of it.
  \item Run GHDL's first two steps on the empty top level, read what each reports, and say what the
    third step requires that does not exist yet.
  \item Say what positional port binding buys and what it gives up.
\end{itemize}

\needspace{6\baselineskip}
\subsection*{Questions to test yourself with}
\begin{enumerate}
  \item Why does \code{bit_timer} own only \emph{when}, and not \emph{what}?
  \item Which block knows the CAN frame's format, and why only one?
  \item Why does \code{can_controller} synchronise its own asynchronous inputs instead of requiring
    whoever instantiates it to do so?
  \item What breaks if you swap two ports of the same type, and when would you notice?
  \item What do \code{ghdl -a}, \code{ghdl -e} and \code{ghdl -r} each do?
  \item Which register-mapped operation corresponds to \code{tx_req}, and which to \code{tx_done}?
\end{enumerate}

\needspace{6\baselineskip}
\subsection*{Target}
The group's repository should exist, with \code{main} protected, and the files handed out should
have reached \code{main} through a reviewed pull request. The top level's entity should exist and analyse
cleanly. After that it is up to the group: someone may well start on \code{bit_timer} or
\code{crc15} straight away. The four leaf modules, \code{bit_timer}, \code{crc15},
\code{tx_shift_reg} and \code{rx_shift_reg}, do not depend on each other and each has a testbench of
its own, so a group of 4--5 can build them in parallel. It is the one point in the project where the
order genuinely matters to how fast you make progress; see appendix~\ref{app:project}.

\needspace{6\baselineskip}
\subsection*{Looking ahead}
\Chapref{c12:ch} builds the first two building blocks going into the empty architecture, both small:
the synchroniser \code{meta_prev}, which knows nothing at all, and \code{bit_timer}, the first block
that knows something about CAN.
